\begin{statement}[C.4.2; Proved: the full division identity]
Let $\psi_m$ be the standard division polynomial of a smooth
Weierstrass model, normalized by $\psi_1=1$, $\psi_2=2y$.
For $m\ge2$ and $Q,mQ\ne O$,
\begin{equation}\label{eq:cs-division}
 N(Q)^{m^2}
       =N(mQ)\frac{|\Delta_E|^{(m^2-1)/12}}{|\psi_m(Q)|}.
\end{equation}
For the displayed family this specializes to
\[
 N(P)^4=3^{3/4}|t-1|^{-1/4}N(2P),\qquad
 N(P)^9=16|t|^{-2/3}N(3P).
\]
If $P_t$ has exact order $n$ and $m\ge2$ satisfies
$m\equiv1$ or $-1\pmod n$, then
\begin{equation}\label{eq:cs-torsionvalue}
 N(P_t)^{12}
       =\frac{|\Delta(t)|}{|\psi_m(P_t)|^{12/(m^2-1)}}.
\end{equation}
The positive value on the right is independent of $m$.
\end{statement}
\begin{proof}
Use the complex uniformization $x=\wp_\Lambda(z)$,
$y=\wp'_\Lambda(z)/2$. The rational function $\psi_m$ has divisor
$\sum_{0\ne Q\in E[m]}(Q)-(m^2-1)(O)$ and leading Laurent term
$(-1)^{m-1}m z^{1-m^2}$. One way to verify both normalizations without
a model convention is to form
$(-1)^{m-1}\sigma_\Lambda(mz)/\sigma_\Lambda(z)^{m^2}$.
Its lattice factors cancel, its zeros are exactly the nonzero
$m$-division points, and its expansion at zero is the displayed
one. The addition formula for $\sigma$ gives the standard
division-polynomial recurrence with initial values $1,2y$, so it is
the stated rational function.

The function
\[
 m^2U_\Lambda(z)-U_\Lambda(mz)
                    +\log|\psi_m(Q(z))|
\]
is harmonic after its logarithmic singularities cancel. Its
Laplacian has zero constant term as well, so it is constant.
The expansion in C.1.1 and the leading term of $\psi_m$ make that
constant $(m^2-1)\log|\Delta_E|/12$. This proves
\eqref{eq:cs-division}. Direct substitution of
\[
 \psi_2(P)=t^2(t-1),\qquad
 \psi_3(P)=-9t^6(t-1)^2/16
\]
gives the two special identities. For $m\equiv\pm1\pmod n$,
$mP=\pm P$ and $N$ is even. Cancelling $N(P)$ in
\eqref{eq:cs-division} gives \eqref{eq:cs-torsionvalue}.
Every right side is the same positive number $N(P)^{12}$, which
proves independence of $m$, not merely numerical agreement.
\end{proof}
